\documentclass[12pt,a4paper]{article}

% -------------------------
% PAQUETES
% -------------------------
\usepackage[utf8]{inputenc}
\usepackage[T1]{fontenc}
\usepackage[spanish]{babel}

\usepackage{amsmath}
\usepackage{amssymb}
\usepackage{geometry}
\usepackage{enumitem}

\geometry{
    margin=2.5cm
}

% -------------------------
% DATOS
% -------------------------
\title{Modelos y Sistemas\\
Primer Parcial}
\author{}
\date{}

% -------------------------
% DOCUMENTO
% -------------------------
\begin{document}

\maketitle

\section*{Ejercicio 1}

Consideramos un sistema masa-resorte-amortiguador con:

\[
m = 7\,kg
\]

\[
c = 2100\,\frac{Ns}{m}
\]

\[
k = 630\,\frac{N}{m}
\]

y condiciones iniciales:

\[
y(0)=1\,m
\]

\[
\dot{y}(0)=0.21\,\frac{m}{s}
\]

\subsection*{a) Ecuación diferencial}

La ecuación del sistema es:

\[
m\ddot{y}(t)+c\dot{y}(t)+ky(t)=u(t)
\]

Reemplazando los valores:

\[
7\ddot{y}(t)+2100\dot{y}(t)+630y(t)=u(t)
\]

\subsection*{b) Representación en espacio de estados}

Definimos las variables de estado:

\[
x_1(t)=y(t)
\]

\[
x_2(t)=\dot{y}(t)
\]

Entonces:

\[
\dot{x}_1=x_2
\]

A partir de la ecuación diferencial:

\[
7\ddot{y}+2100\dot{y}+630y=u
\]

obtenemos:

\[
\ddot{y}
=
-\frac{2100}{7}\dot{y}
-\frac{630}{7}y
+\frac{1}{7}u
\]

por lo tanto:

\[
\dot{x}_2
=
-300x_2-90x_1+\frac{1}{7}u
\]

La representación matricial queda:

\[
\begin{bmatrix}
\dot{x}_1\\
\dot{x}_2
\end{bmatrix}
=
\begin{bmatrix}
0 & 1\\
-90 & -300
\end{bmatrix}
\begin{bmatrix}
x_1\\
x_2
\end{bmatrix}
+
\begin{bmatrix}
0\\
\frac{1}{7}
\end{bmatrix}u
\]

\subsection*{c) Polinomio característico y solución general}

El polinomio característico se obtiene de:

\[
\det(sI-A)=0
\]

con

\[
A=
\begin{bmatrix}
0 & 1\\
-90 & -300
\end{bmatrix}
\]

Entonces:

\[
\begin{vmatrix}
s & -1\\
90 & s+300
\end{vmatrix}
=0
\]

por lo que:

\[
s(s+300)+90=0
\]

\[
s^2+300s+90=0
\]

Las raíces son:

\[
s_{1,2}
=
\frac{-300\pm\sqrt{300^2-4(90)}}{2}
\]

Por lo tanto, la solución general es de la forma:

\[
y(t)=C_1e^{s_1t}+C_2e^{s_2t}
\]

Las constantes se determinan utilizando las condiciones iniciales.

\subsection*{d) Matriz de transición de estados}

La matriz de transición de estados se obtiene como:

\[
\Phi(t)=e^{At}
\]

donde $A$ es la matriz del sistema.

También puede calcularse mediante la transformada de Laplace:

\[
\Phi(t)
=
\mathcal{L}^{-1}
\left\{
(sI-A)^{-1}
\right\}
\]

\newpage

% =====================================================
% EJERCICIO 2
% =====================================================

\section*{Ejercicio 2}

Se tiene el sistema:

\[
\ddot{y}(t)+8\dot{y}(t)+7y(t)=u(t)
\]

\subsection*{a) Función transferencia}

Aplicando transformada de Laplace y considerando condiciones iniciales nulas:

\[
s^2Y(s)+8sY(s)+7Y(s)=U(s)
\]

Factorizando:

\[
Y(s)(s^2+8s+7)=U(s)
\]

Por definición:

\[
H(s)=\frac{Y(s)}{U(s)}
\]

entonces:

\[
\boxed{
H(s)=\frac{1}{s^2+8s+7}
}
\]

Factorizando el denominador:

\[
s^2+8s+7=(s+1)(s+7)
\]

por lo tanto:

\[
\boxed{
H(s)=\frac{1}{(s+1)(s+7)}
}
\]

\subsection*{b) Respuesta al impulso}

Para una entrada impulso:

\[
u(t)=\delta(t)
\]

su transformada es:

\[
U(s)=1
\]

Entonces:

\[
Y(s)=H(s)
\]

por lo tanto:

\[
Y(s)=\frac{1}{(s+1)(s+7)}
\]

Realizamos fracciones simples:

\[
\frac{1}{(s+1)(s+7)}
=
\frac{A}{s+1}+\frac{B}{s+7}
\]

De donde:

\[
A=\frac{1}{6},
\qquad
B=-\frac{1}{6}
\]

Entonces:

\[
Y(s)
=
\frac{1}{6}\frac{1}{s+1}
-
\frac{1}{6}\frac{1}{s+7}
\]

Aplicando transformada inversa:

\[
\boxed{
h(t)=
\frac{1}{6}e^{-t}
-
\frac{1}{6}e^{-7t}
}
\]

\subsection*{c) Respuesta al escalón}

Para una entrada escalón:

\[
u(t)=1(t)
\]

su transformada es:

\[
U(s)=\frac{1}{s}
\]

Entonces:

\[
Y(s)=H(s)U(s)
\]

\[
Y(s)
=
\frac{1}{s(s+1)(s+7)}
\]

Realizamos fracciones simples:

\[
\frac{1}{s(s+1)(s+7)}
=
\frac{A}{s}
+
\frac{B}{s+1}
+
\frac{C}{s+7}
\]

Obtenemos:

\[
A=\frac{1}{7},
\qquad
B=-\frac{1}{6},
\qquad
C=\frac{1}{42}
\]

Por lo tanto:

\[
Y(s)
=
\frac{1}{7s}
-
\frac{1}{6(s+1)}
+
\frac{1}{42(s+7)}
\]

Aplicando transformada inversa:

\[
\boxed{
y(t)
=
\frac{1}{7}
-
\frac{1}{6}e^{-t}
+
\frac{1}{42}e^{-7t}
}
\]

\newpage

% =====================================================
% CONCEPTOS TEÓRICOS
% =====================================================

\section*{Conceptos importantes}

\subsection*{Función transferencia}

La función transferencia se define como:

\[
\boxed{
H(s)=\frac{Y(s)}{X(s)}
}
\]

considerando condiciones iniciales nulas.

\subsection*{Respuesta al impulso}

Si la entrada es:

\[
x(t)=\delta(t)
\]

entonces:

\[
X(s)=1
\]

y por lo tanto:

\[
Y(s)=H(s)
\]

La respuesta al impulso es:

\[
\boxed{
h(t)=\mathcal{L}^{-1}\{H(s)\}
}
\]

\subsection*{Convolución}

Para un sistema LTI:

\[
\boxed{
y(t)=
\int_{-\infty}^{\infty}
x(\sigma)h(t-\sigma)\,d\sigma
}
\]

En Laplace:

\[
\boxed{
Y(s)=H(s)X(s)
}
\]

\subsection*{Realimentación negativa}

Para el esquema general:

\[
\boxed{
\frac{Y(s)}{X(s)}
=
\frac{G_1}{1+G_1G_2}
}
\]

\subsection*{Circuito RC pasabajo}

La función transferencia es:

\[
\boxed{
H(s)=\frac{1}{1+sRC}
}
\]

\subsection*{Teorema del valor inicial}

\[
\boxed{
y(0)=
\lim_{s\rightarrow\infty}sY(s)
}
\]

\subsection*{Teorema del valor final}

\[
\boxed{
y(\infty)=
\lim_{s\rightarrow0}sY(s)
}
\]

siempre que se cumplan las condiciones de estabilidad.

\subsection*{Entradas típicas}

Impulso:

\[
\delta(t)
\quad\longleftrightarrow\quad
1
\]

Escalón:

\[
1(t)
\quad\longleftrightarrow\quad
\frac{1}{s}
\]

Rampa:

\[
t
\quad\longleftrightarrow\quad
\frac{1}{s^2}
\]

Sinusoidal:

\[
\sin(\omega_0t)
\quad\longleftrightarrow\quad
\frac{\omega_0}{s^2+\omega_0^2}
\]

Exponencial:

\[
e^{at}
\quad\longleftrightarrow\quad
\frac{1}{s-a}
\]

\subsection*{Diagrama de Bode}

La magnitud de la función transferencia en decibeles es:

\[
\boxed{
|H(j\omega)|_{dB}
=
20\log |H(j\omega)|
}
\]

Los parámetros importantes son:

\begin{itemize}
    \item Magnitud.
    \item Fase.
    \item Frecuencia de corte.
    \item Ancho de banda.
    \item Resonancia.
\end{itemize}

\end{document}
